% !TEX root = ./ECEN-635 - Homework 2 - Brody Jordan.tex

\documentclass[11pt]{article}

\usepackage[
    assignmentDueDate={September 10, 2025 at 4:30 pm}
]{C:/Users/bljor/Sync/Courses/assignment}

\usepackage{siunitx}

\begin{document}
\makeAssignmentTitle

% --- Problem 1 ---
\begin{homeworkProblem}

Convert the following time domain expressions for the electric and magnetic field intro frequency and domain expressions. 
Show each mathematical step in your calculations. A miraculous answer without the complete calculation will be judged harshly.

\begin{enumerate}[label=(\alph*), topsep=5pt, itemsep=0pt]
    \item $\mathbf{E}(t)=4 \cos (\omega t-28 z) \hat{\mathbf{y}} \;\; \unit{V/m}$
    \item $\mathbf{E}(t)=9 \sin (\omega t-23 z) \hat{\mathbf{x}} \;\; \unit{V/m}$
    \item $\mathbf{H}(t)=7 \cos (\omega t-23 z) \hat{\mathbf{x}} \;\; \unit{mA/m}$
    \item $\mathbf{E}(t)=4.6 e^{-48 \pi z} \cos (\omega t-139 z+\pi / 4) \hat{\mathbf{y}} \;\; \unit{V/m}$
\end{enumerate} \vspace{1\baselineskip}

\solution



\end{homeworkProblem}

\pagebreak

% --- Problem 2 ---
\begin{homeworkProblem}

Convert the following frequency domain expressions to the time domain. Show each mathematical step in your calculation.

\begin{enumerate}[label=(\alph*), topsep=5pt, itemsep=0pt]
    \item $\mathbf{E} = 15 e^{-j 23 z} \hat{\mathbf{x}} \;\; \unit{V/m}$
    \item $\mathbf{E} = 32 e^{-(52 + j 135) z} \hat{\mathbf{y}} \;\; \unit{V/m}$
    \item $\mathbf{H} = [(2 + j2) \hat{\mathbf{x}} + (1 - j 3) \hat{\mathbf{y}}] e^{-j 245 z} \;\; \unit{A/m}$
    \item $\rho = 3 j e^{j\pi/2} \;\; \unit{C/m^3}$ \\
    Be sure to place each of the vector components in simplest magnitude and phase form before
    converting to the time domain.
\end{enumerate} \vspace{1\baselineskip}

\solution



\end{homeworkProblem}

\pagebreak

% --- Problem 3 ---
\begin{homeworkProblem}

Assuming a homogeneous region where there is no current, charge, or loss:
\begin{enumerate}[label=(\alph*), topsep=5pt, itemsep=0pt]
    \item Use Maxwell's equations to convert the frequency domain electric field expression in 2(b) to the corresponding
    magnetic field. Show all mathematical steps.
    \item Use Maxwell's equations to convert the frequency domain magnetic field in 2(c) into a frequency domain electric
    field into a frequency domain magnetic field. Show all mathematical steps.
\end{enumerate} \vspace{1\baselineskip}

\solution

\end{homeworkProblem}

\pagebreak

% --- Problem 4 ---
\begin{homeworkProblem}

Use 1(b) and the corresponding electric field found in problem 1(c) to calculate the time average power density. 
Show all mathematical steps. \vspace{1\baselineskip}

\solution


\end{homeworkProblem}

\pagebreak

% --- Problem 5 ---
\begin{homeworkProblem}

Starting with the frequency domain form of Maxwell's equations, derive the magnetic field vector Helmholtz equation. Number
your Maxwell equations and at every step write briefly which Maxwell equation and/or which identity you are using
from the Vectory Identities section of Appendix II.3.1 (pg. 958). Assume there is no electric or magnetic current and
no electric or magnetic charge. \vspace{1\baselineskip}

\solution



\end{homeworkProblem}

\end{document}